% supp_front.tex -- heading of the supplementary material (second part of the same PDF) and
% its numbering: appendices S1, S2, ..., statements S1.1, ..., tables S1, S2, ...; equations continue.
\clearpage
\makeatletter
\setcounter{section}{0}
\renewcommand{\thesection}{S\arabic{section}}
\@ifundefined{theHsection}{}{\renewcommand{\theHsection}{S\arabic{section}}}
\setcounter{table}{0}
\renewcommand{\thetable}{S\arabic{table}}
\@ifundefined{theHtable}{}{\renewcommand{\theHtable}{S\arabic{table}}}
\makeatother
\phantomsection\label{part:supp}
\addcontentsline{toc}{part}{Supplementary material}
\begin{center}
{\Large\bfseries Supplementary material}\par\smallskip
to ``The $6\times6\times6$ tic-tac-toe game is a draw''
\end{center}
\bigskip

\noindent
The appendices S1 to S7 belong to the same paper but not to its main line. (The programs and the data
are not here but in the supplementary archive described in \cref{sec:repro}.) Nothing in them is used in the
proof of \cref{thm:main} or of \cref{cor:seven}, and nothing in them is needed in order to replicate the
verification of \cref{fact:H}: the definitions, the two verification algorithms with their correctness
proofs, the file formats, the canonical form, the test vectors and the numbers of states in every round
are contained in \cref{sec:board,sec:pairing,sec:leaf,sec:game,sec:certificate,sec:verification} and in
\cref{sec:formats}. The supplementary appendices contain results that the main theorem does not use, the
description of the search that found the certificate, the step-by-step description of the leaf test of
the main verifier, the details of the audits and of the cross-checks, the integer program, the details of the archive and of the reproduction of the computation, and tables
that repeat information of the main part in another arrangement; \cref{tab:supp-overview} gives an
overview. The
convention on Computed Facts of \cref{ssec:intro-nature} applies here as well: statements called Theorem,
Proposition, Lemma or Corollary are proved, and statements established by running programs are called
Computed Facts or are reported with their source.

\begin{table}[ht]
\centering\small
\begin{tabular}{@{}l>{\raggedright\arraybackslash}p{0.56\textwidth}l@{}}
\toprule
appendix & content & supplements\\
\midrule
\ref{sec:extra} & the full automorphism group of the board (Computed Fact); the leaf test in words;
   monotonicity of the pairing condition in \Breaker's stones & \cref{sec:board,sec:leaf}\\
\addlinespace[2pt]
\ref{sec:extra-game} & lazy monotonicity and the eager variant of the abstract game; the exact numbers of
   stones of the patterns $\Etwo$ and $\Ethree$ & \cref{sec:game}\\
\addlinespace[2pt]
\ref{sec:supp-search} & how the search works; two observations made with the search program alone &
   \cref{ssec:search,ssec:observations}\\
\addlinespace[2pt]
\ref{sec:supp-checks} & the leaf test of the main verifier step by step, with its statistics; mutation
   experiments and comparisons with brute force; exact enumeration, integer program and simulated games;
   known minor issues of the tools & \cref{ssec:algleaf,ssec:audits,ssec:cross,ssec:trust}\\
\addlinespace[2pt]
\ref{sec:ip} & the integer program used as a cross-check & \cref{ssec:cross}\\
\addlinespace[2pt]
\ref{sec:supp-repro} & contents of the archive; outputs and running times of the verifiers; the steps of
   the reproduction script & \cref{sec:repro}\\
\addlinespace[2pt]
\ref{sec:supp-tables} & the geometry in terms of cell ids; the numbers of states in every round, by rows;
   orbit counts; the minimum number of stones for a pairing & \cref{sec:formats}\\
\bottomrule
\end{tabular}
\caption{The supplementary appendices.}
\label{tab:supp-overview}
\end{table}
